% TODO:
%
%
% - Add the trolly problem as an introductory case.
% - Discuss exploitability vs. oversight on a controllability spectrum. Maybe your own paper by now. 
% - Maybe rather use autonomous driving as a case study, since it is more relatable and has more real-world examples of accountability gaps. The trolley problem is a bit too abstract and hypothetical, and the recruiting tool is a bit too niche and technical. Autonomous driving has both high stakes and real cases of accidents, recalls, and liability disputes.
%




\chapter{Accountability}\index{Accountability}
\printchapterversion{02_Accountability}

\newthought{The Responsibility Gap}

\marginnote{%
    \begin{flushright}
    \newthought{In der Strafkolonie}\\
    Franz Kafka (1919)\\
    {\color{black}\qrcode[height=0.9cm]{https://de.wikisource.org/wiki/In_der_Strafkolonie}}
    \end{flushright}%
    \medskip
    \newthought{Minority Report}\cite{spielberg2002minority} A
    pre-crime unit arrests people before they offend, based on
    algorithmic prediction. When the system targets its own chief, the
    question of who holds the algorithm accountable becomes impossible
    to avoid.%
}

% As a leading group, you are required to moderate this session. Add margin notes in this script in case you want to add any additional questions or points to the discussion.

\newthought{A visiting traveller} witnesses a penal machine execute a condemned man for a
sentence no living person can read or explain, and is invited to endorse the system.\\

\newthought{Algorithmic governance} and black-box AI in high-stakes domains highlight
opacity and accountability gaps: audits, explainability requirements, and governance mechanisms
all attempt to close the distance between automated decision and human answerability.

\vspace{3.5em}

\begin{tabular}{p{3.8cm} p{6.5cm}}
    \textbf{Required reading} &
        \textit{In der Strafkolonie}\cite{kafka1919strafkolonie} by Franz Kafka
        (1919; Reclam UB~9900) \\[3pt]
    \textbf{Preparation}      &
        Complete the Guided Reading Sheet individually before the session. \\
\end{tabular}

\vspace{3.5em}

\section*{Guided Reading Sheet}

\noindent Read \emph{In der Strafkolonie} before the session. Work through the
questions below individually, then bring your notes to the discussion.

\begin{enumerate}
    \item The condemned man never learns what he is accused of or what
    the machine will inscribe on his body. What does this tell us about
    the relationship between legibility and legitimacy? Can a system be
    just if it cannot explain itself to those it acts upon?

    \item Three figures surround the machine: the officer who
    understands and maintains it, the condemned man who cannot, and
    the visiting traveller who is invited to endorse it. How does
    responsibility shift across this triangle? Who, if anyone, is
    morally culpable for what happens?

    \item The story illustrates the idea of the \emph{responsibility
    gap}: when a complex system causes harm, no single human actor is
    clearly at fault. Can you name a contemporary case (in law,
    medicine, finance, or infrastructure) where this gap has appeared?

    \item The EU AI Act (2024) requires human oversight for high-risk
    AI systems. What does "oversight" mean in practice if the system's
    decisions are too fast, too complex, or too numerous for humans to
    review individually?

    \item Compare \emph{causal responsibility} (who caused the harm)
    with \emph{moral responsibility} (who is blameworthy) and
    \emph{legal liability} (who must compensate). Do these three always
    converge? Give an example where they do not.\\
\end{enumerate}



\section*{Opinion Landscape and Debate}

\noindent \textbf{Format:} Surrounded-style debate. See
Appendix~A for the full protocol, speaker's list rules, and
moderator scripts.

\medskip
\noindent \textbf{Opening question:} If a system causes injustice at this
scale, and every individual who used it acted within their authorised role,
is anyone guilty of wrongdoing; and if not, what does that mean for the
concept of accountability?
\newthought{Opinion~A, Accountability.}\quad Named human accountability
is achievable and required: when an autonomous system causes harm, the
architecture of governance must ensure that specific individuals remain
legally and morally answerable for its outputs.
\vspace{1.5em}


\medskip
\noindent\textbf{The Defence Lawyer} argues the tool violated the
    defendant's right to a fair trial and demands a named responsible party.
    \begin{itemize}
        \item \textit{My client cannot cross-examine a dataset. She cannot
        see the features the model weighted. She cannot confront the humans
        whose historical decisions trained it. That is not a fair trial;
        it is a verdict delivered by an oracle.}
        \item \textit{The right to a reasoned decision is foundational to
        due process. The algorithm said so is not a reason. It is an
        abdication.}
        \item \textit{If a biased human expert testified in court, we could
        challenge their credentials, their methods, their prior record. We
        have no equivalent procedure for an algorithm. That asymmetry is a
        constitutional problem.}
    \end{itemize}

\medskip
\noindent\textbf{The Regulator} proposes a new accountability
    framework and must defend it against both sides.
    \begin{itemize}
        \item \textit{The existing categories (product liability,
        professional negligence, criminal intent) were designed for
        individual actors and single causal chains. None of them fit
        distributed, automated decision systems. We need a new category:
        algorithmic liability.}
        \item \textit{Auditability must be a pre-condition for deployment in
        high-stakes domains, not a post-incident investigation. If you cannot
        explain why your system made a decision before harm occurs, you
        should not be permitted to deploy it.}
        \item \textit{The question is not who is liable for this specific
        incident. It is what governance regime prevents the next thousand
        incidents at scale.}
    \end{itemize}



\newthought{Opinion~B, Compliance.}\quad The existing system already
provides accountable humans; the motion is either already met or
unworkable: demanding named individual accountability for each algorithmic
output ignores how distributed decision-making necessarily operates.
\vspace{1.5em}


\medskip
\noindent\textbf{The Judge} defends reliance on the tool as
    consistent, evidence-based sentencing practice.
    \begin{itemize}
        \item \textit{Documented bias in human sentencing (based on time
        of day, defendant ethnicity, and courtroom demeanour) is
        well-established in the research literature. The tool reduces that
        variance. Less variance is more equal treatment.}
        \item \textit{I did not delegate my judgement to the tool. I
        consulted it as one input among many, as I would a psychiatric
        assessment or a victim impact statement. The tool is not the
        verdict.}
        \item \textit{If we ban algorithmic tools from courtrooms, we return
        to a system where the most important variable in a sentencing outcome
        is which judge you happened to draw. Is that preferable?}
    \end{itemize}

\medskip
\noindent\textbf{The Software Vendor} insists the tool performs
    within documented parameters and liability rests with the deployer.
    \begin{itemize}
        \item \textit{Our documentation states in the executive summary that
        human review is mandatory before any recommendation is acted upon.
        The court chose to rely on the output without applying that review.
        We cannot be held liable for misuse.}
        \item \textit{The training data was provided by the court system
        itself. If the historical sentencing record reflects bias, that is a
        problem with the justice system, not with a model that learned from
        it.}
        \item \textit{Product liability law requires that a product fails to
        perform as documented. Ours did not. Extending vendor liability to
        downstream misuse would make the development of any decision-support
        tool legally untenable.}
    \end{itemize}


\newthought{Key Learnings}
The responsibility gap, the structural absence of a single culpable
actor in distributed sociotechnical systems, is not caused by
negligence but by the architecture of those systems: standard
categories of fault (criminal intent, professional negligence, product
liability) apply only awkwardly when causation is diffuse. Closing it
requires deliberate governance design, not post-hoc blame. Explainability, the
property of being able to account for outputs in terms a human
decision-maker can evaluate, is a necessary but not sufficient
condition for accountability; a system can be interpretable and still
have no human willing to own its outputs. Meaningful human control, a
regulatory concept that requires humans to retain the ability to
understand, intervene in, and reverse AI decisions, is contested in
practice: it may mean the ability to intervene, to understand, to
reverse, or merely to authorise, and these are different requirements
with different institutional costs. The lessons of aviation safety culture, medical liability law,
and financial auditing all offer partial models for algorithmic governance, but none transfers cleanly.

\medskip
\noindent\textit{Teaser: accountability assumes we can separate the
system from its surroundings. But what happens when the AI is no
longer a tool we hold, but a body that acts in the world alongside
us?}

\clearpage
